\part{El espectro como recurso legal}
\label{part:regulacion}

\chapter{Por qué se regula el espectro}

\section{Qué es}

El espectro radioeléctrico es un recurso \textbf{compartido y finito}: dos transmisores en la
misma frecuencia, al mismo tiempo y en el mismo lugar, se interfieren mutuamente y ninguno se
puede decodificar. A diferencia de un cable, donde añadir más usuarios es cuestión de tender más
cobre, el espectro no se puede ``fabricar'': todo el que existe ya está ahí, y su uso debe
coordinarse por ley para que servicios críticos —aviación, emergencias, telefonía celular
licenciada— no sufran interferencia de dispositivos no autorizados.

Por eso cada país tiene una autoridad reguladora que decide qué franjas de frecuencia se
reservan para qué uso, y bajo qué condiciones técnicas (potencia máxima, ancho de banda, ciclo
de trabajo) se puede transmitir en cada una.

\section{Qué son las bandas ISM}

Las bandas \textbf{ISM} (\emph{Industrial, Scientific and Medical}) son franjas de espectro
que la mayoría de países reserva para uso \textbf{libre y sin licencia individual}, siempre que
el equipo cumpla ciertas condiciones técnicas (potencia máxima, tipo de modulación). Es el
régimen bajo el cual operan wifi (\SI{2.4}{GHz}), Bluetooth, y también las bandas sub-GHz que
usa LoRa. ``Libre'' no significa ``sin reglas'': significa que no hace falta pedir una licencia
de operación individual, pero el equipo sí debe respetar los límites técnicos de la banda.

\section{Por qué importa aquí}

El proyecto \texttt{lorasync} depende enteramente de que su banda de operación sea, en efecto,
de uso libre en Colombia. Transmitir en una banda licenciada sin autorización —aunque sea con
poca potencia— es una infracción, no un simple detalle técnico.

\chapter{Qué es la ANE}

La \textbf{Agencia Nacional del Espectro} (ANE) es la entidad colombiana que administra,
vigila y controla el uso del espectro radioeléctrico en el país. Publica el \textbf{Cuadro
Nacional de Atribución de Bandas de Frecuencias} (CNABF), el documento que determina, banda por
banda, qué servicios están autorizados y bajo qué condiciones. Es la referencia oficial usada
para las cifras de esta sección~\cite{ane_cnabf}.

\chapter{Por qué 915--928 MHz sí y 868 MHz no}

\section{Qué es}

En Europa, la banda sub-GHz de uso libre para IoT es \SI{863}{}--\SI{870}{MHz}, con
\SI{868}{MHz} como frecuencia de referencia habitual —de ahí que muchísima documentación y
hardware genérico de LoRa (incluidas antenas) esté etiquetado ``868 MHz''—. Colombia
\textbf{no} sigue la atribución europea en esa franja: la ANE clasifica
\SI{851}{}--\SI{915}{MHz} como espectro celular licenciado (usado por los operadores móviles),
y reserva la banda de uso libre en \SI{915}{}--\SI{928}{MHz}, alineada con el modelo
norteamericano de banda ISM de 900 MHz~\cite{ttn_au915}.

\begin{figure}[h]
  \centering
  \begin{tikzpicture}[scale=1.0, font=\footnotesize, every node/.style={align=center}]
    % eje
    \draw[-{Latex}] (0,0) -- (10.5,0) node[right, font=\footnotesize] {frecuencia (MHz)};
    % banda restringida (celular licenciado)
    \fill[red!25] (0.5,0.3) rectangle (7,1.7);
    \draw (0.5,0.3) rectangle (7,1.7);
    \node[text width=5.6cm] at (3.75,1.0) {celular licenciado\\(851--915 MHz)\\\textbf{prohibido}};
    % banda libre
    \fill[green!25] (7,0.3) rectangle (9,1.7);
    \draw (7,0.3) rectangle (9,1.7);
    \node at (8,1.0) {\textbf{915--928}\\libre};
    % ticks y etiquetas del eje
    \foreach \x/\lbl in {0.5/851, 7/915, 9/928} {
      \draw (\x,0.25) -- (\x,0.35);
      \node[below, font=\scriptsize] at (\x,0.15) {\lbl};
    }
    % marca 868 (bajo el eje, separada de las etiquetas de los ticks)
    \draw[dashed, blue!60!black, thick] (4.55,-0.9) -- (4.55,1.9);
    \node[blue!60!black, font=\scriptsize, text width=3.2cm] at (4.55,-1.15)
      {868 MHz (Europa)\\\textbf{no usar en CO}};
    % marca canal 70 -> 920 MHz (sobre el eje)
    \draw[dashed, orange!70!black, thick] (7.9,-0.9) -- (7.9,2.1);
    \node[orange!70!black, font=\scriptsize, text width=3.2cm] at (7.9,2.35)
      {canal 70 = 920 MHz\\(usado aquí)};
  \end{tikzpicture}
  \caption{La banda de \SI{868}{MHz}, libre en Europa, cae dentro del espectro celular
    licenciado en Colombia. La banda libre colombiana equivalente es 915--928\,MHz.}
  \label{fig:banda-colombia}
\end{figure}

\section{Por qué importa aquí}

Confundir la convención europea con la colombiana es el error regulatorio más probable en un
proyecto que usa hardware genérico de LoRa (documentación, foros y valores de fábrica casi
siempre asumen \SI{868}{MHz}). El valor de fábrica del módulo Waveshare (canal 18) corresponde
precisamente a \SI{868}{MHz}, y \textbf{no debe usarse en Colombia}.

\section{La fórmula de canalización}

El firmware del módulo no acepta una frecuencia directamente: acepta un número de
\textbf{canal}, que se traduce a frecuencia con:

\begin{equation}
  f_{\text{MHz}} = 850 + \text{canal}
  \label{eq:canal}
\end{equation}

Los canales válidos para la banda libre colombiana son, por tanto, los que caen entre 915 y
928: $\text{canal} \in [65, 78]$. El proyecto usa el \textbf{canal 70}, que por la
ecuación~\eqref{eq:canal} da $850 + 70 = \SI{920}{MHz}$, cómodamente dentro del rango legal y
lejos de sus bordes.

\begin{trampa}
El valor de fábrica del módulo es el \textbf{canal 18}, que da $850 + 18 = \SI{868}{MHz}$: la
frecuencia europea, ilegal en Colombia. Cualquier módulo recién desempacado o reseteado a
valores de fábrica transmite fuera de banda hasta que se reconfigura explícitamente.
\end{trampa}

\chapter{Ciclo de trabajo y convivencia}

\section{Qué es}

El \textbf{ciclo de trabajo} (\emph{duty cycle}) es el porcentaje de tiempo que un transmisor
tiene permitido ocupar el canal. En Europa, la banda de \SI{868}{MHz} impone un límite estricto
de ciclo de trabajo (típicamente 1\,\% por sub-banda) como mecanismo regulatorio para forzar a
que muchos dispositivos puedan compartir el mismo canal sin coordinación explícita. El modelo
norteamericano (FCC Part 15.247), del que deriva la banda libre de 900 MHz en Colombia, no
impone un límite de ciclo de trabajo tan estricto, pero exige en su lugar \textbf{ancho de
banda mínimo} o \textbf{salto de frecuencia} (\emph{frequency hopping}), como mecanismo
alternativo para evitar que un solo dispositivo monopolice una franja estrecha del
espectro~\cite{fcc_15247}.

\section{Por qué importa aquí}

Los módulos Waveshare de este proyecto son de \textbf{canal fijo}: no implementan salto de
frecuencia. Esto significa que no pueden cumplir la variante ``frequency hopping'' del modelo
FCC. La variante que sí aplica es la de \textbf{ancho de banda mínimo mediante modulación
digital de banda ancha}: los \SI{500}{kHz} de ancho de banda elegidos para este proyecto no son
solo la opción de máxima capacidad (menor tiempo de aire, sección~\ref{sec:airtime}) —también
son la que mejor encaja con esta condición regulatoria—.

\begin{nota}
Esta interpretación es razonable para un prototipo de investigación con 2--30 nodos, pero
\textbf{no sustituye una verificación formal con la ANE} antes de cualquier despliegue
permanente o comercial. Es uno de los supuestos declarados explícitamente en el plan del
proyecto.
\end{nota}

Con solo 2--30 nodos transmitiendo tramas cortas de forma coordinada (ver el diseño TDMA de la
Fase~4), la ocupación real del canal —41,5\,\% en el escenario de 30 nodos— queda por debajo de
niveles que pondrían en riesgo la convivencia con otros usuarios de la misma banda libre.

\section{Dónde aparece en el código}

La validación del rango de canal (65--78) vive en \texttt{common/config.py} (Fase~1): el
sistema \textbf{se niega a arrancar} fuera de ese rango salvo que se pase explícitamente el
flag \texttt{allow\_out\_of\_band}, precisamente para que un error de configuración no termine
transmitiendo en \SI{868}{MHz} por accidente.
